Enhancers regulate gene expression in a cell-type- and context-specific manner, and linking them to their target genes is essential for interpreting disease-associated noncoding variants. Single-cell RNA and ATAC sequencing allow chromatin accessibility peaks to be linked to genes within individual cell types, but existing methods rely on the correlation between peak accessibility and gene expression across cells or across individuals, and correlation does not imply causation. Here we present a Mendelian randomization (MR) framework that uses chromatin accessibility and expression quantitative trait loci from population-scale single-cell data to estimate the causal effect of peak accessibility on gene expression within a cell type. The framework combines inverse-variance weighted, MR-Egger and summary-data-based MR estimators with a joint test, accommodating correlated instruments and allowing for directional pleiotropy. We applied the framework to naive B cells, naive CD4$^+$ T cells and naive CD8$^+$ T cells from the Chinese Immune Multi-Omics Atlas and compared it with two correlation-based methods. MR identified fewer peak--gene links than either correlation-based method, but some links were detected only by MR. For links detected by all methods, the MR estimates almost always agreed with one another in direction, whereas disagreement arose mainly with the correlation across metacells. MR-supported links generally had higher odds ratios of overlap with enhancer-driven regulons from the same atlas and were enriched in six external regulatory references. Peaks of links detected only by MR were 3.28 times as likely as candidate peaks to overlap an H3K27ac HiChIP contact with the same gene and had higher point estimates of heritability enrichment for seven immune-mediated diseases. At disease-associated loci, these links suggested candidate target genes and cell types, including \textit{ERAP1}, \textit{ERBB3} and \textit{THBS3}. Our results show that MR with population-scale single-cell data complements correlation-based mapping of enhancers to target genes.
